\documentclass[../../main]{subfiles}

\begin{document}

\section{Famílias de Conjuntos}
\label{sec:matematica:teoria-conjuntos:familias-de-conjuntos}

\begin{defn}[Família de Conjuntos]
	\label{def:matematica:teoria-conjuntos:familias-de-conjuntos:familia}

	Uma família de conjuntos é uma coleção de conjuntos identificados por
	elementos de um conjunto de índices \(I\).

	Uma família indexada por \(I\) é denotada por

	\[
		\{A_i\}_{i \in I}.
	\]

	Para cada índice \(i \in I\), associa-se um conjunto \(A_i\).

\end{defn}

\begin{defn}[Conjunto de Índices]
	\label{def:matematica:teoria-conjuntos:familias-de-conjuntos:indices}

	O conjunto \(I\) utilizado para identificar os conjuntos de uma
	família é denominado conjunto de índices.
\end{defn}

\begin{obs}

	Também é comum utilizar a notação

	\[
		(A_i)_{i \in I}.
	\]

\end{obs}

\begin{ex}

	A coleção

	\[
		\{A_n\}_{n\in\mathbb{N}}
	\]

	com

	\[
		A_n=\{1,2,\dots,n\}
	\]

	é uma família infinita de conjuntos.

\end{ex}

\begin{defn}[Família Finita]
	\label{def:matematica:teoria-conjuntos:familias-de-conjuntos:familia-finita}

	Uma família

	\[
		\{A_i\}_{i\in I}
	\]

	é denominada finita quando o conjunto de índices \(I\) é finito.

\end{defn}

\begin{defn}[Família Infinita]
	\label{def:matematica:teoria-conjuntos:familias-de-conjuntos:familia-infinita}

	Uma família

	\[
		\{A_i\}_{i\in I}
	\]

	é denominada infinita quando o conjunto de índices \(I\) é infinito.

\end{defn}

\begin{defn}[União Arbitrária]
	\label{def:matematica:teoria-conjuntos:familias-de-conjuntos:uniao-arbitraria}

	Seja

	\[
		\{A_i\}_{i\in I}
	\]

	uma família de conjuntos.

	Define-se

	\[
		\bigcup_{i\in I}A_i
		=
		\{
		x
		\mid
		(\exists i\in I)
		(x\in A_i)
		\}.
	\]

\end{defn}

\begin{defn}[Interseção Arbitrária]
	\label{def:matematica:teoria-conjuntos:familias-de-conjuntos:intersecao-arbitraria}

	Seja

	\[
		\{A_i\}_{i\in I}
	\]

	uma família de conjuntos.

	Define-se

	\[
		\bigcap_{i\in I}A_i
		=
		\{
		x
		\mid
		(\forall i\in I)
		(x\in A_i)
		\}.
	\]

\end{defn}
